\section{AI assignment of reaction direction}\label{sec:s4}

The assignment of a reaction direction from thermodynamics means that nearly half the database doesn't receive a direction ultimately due to the coverage of structures in the database, however well understood the reaction may be. In order to provide an alternative perspective of the reaction direction and to improve the coverage across more reactions in the database we implemented a complementary approach that predicts the direction of the reaction using an ensemble of large language models (LLMs). A prompt was formed that presented only the reaction identifier, its name, and its equation, with the equation written without a direction. None of the data generated from any of the heuristics or the experimental data from openTECR was used in the prompt.

\subsection{Ensemble} Eleven candidate LLMs were scored across the role of proposer, auditor, and adjudicator (see next paragraph on how these are used). They were evaluated based on the rate at which an assigned reaction matched experiments. In the end, we selected: three separate LLMs from different commercial providers as predictors: Claude Opus 4.8, Gemini 3.5 Flash, and GPT-5.6; GPT-5.5 as the auditor and Claude Opus 5 as the adjudicator. The proposers were instructed to respond with a predicted direction ($>$, $<$, $=$, or $?$ where it judges the evidence insufficient), a confidence, and a written justification grounded in the stoichiometry or in known enzyme biology. The auditor compiled a vote table of the proposals and the adjudicator makes the final call.

\subsection{Reactions and status} We filtered out reactions for which the prediction would not have been appropriate, this includes reactions flagged as obsolete, symmetrical transport reactions, and pseudo-reactions that acted as lumped reactions, aggregating many species and for which there is no direction to infer, leaving $\sim46,000$ reactions for which direction is assigned. We store the output of the ensemble in our repository for review, including any objections raised by the auditor.

\subsection{Summary of results} As the set of reactions covers almost all those for which we generated thermodynamics data, we're able to draw a direct comparison. The ensemble commits to a direction far more readily than the thermodynamic rules do, abstaining on $\sim2\%$ of reactions. The confidence that the ensemble returns does not appear to distinguish correct calls when compared to the thermodynamics, and is not appropriate to use as a downstream filter when integrating reactions. Nevertheless, with the far wider coverage, researchers may find the results useful.